% unidades/09-adjetivos.tex
\chapter{🎨 形容詞 — Adjetivos}
\unidadheader{09}{形容詞}{Adjetivos}{Describir cosas, personas y lugares}

\begin{tcolorbox}[vocabbox, title={📌 Vocabulario Nuevo}]
\begin{tabularx}{\linewidth}{llX}
\toprule
\textbf{Japonés} & \textbf{Español} & \textbf{Tipo} \\
\midrule
\vocabitem{きれい}{—}{hermoso / limpio}{adj-な}
\vocabitem{静か}{しずか}{silencioso / tranquilo}{adj-な}
\vocabitem{賑やか}{にぎやか}{concurrido / animado}{adj-な}
\vocabitem{便利}{べんり}{conveniente}{adj-な}
\vocabitem{暇}{ひま}{libre / sin nada que hacer}{adj-na}
\vocabitem{忙しい}{いそがしい}{ocupado}{adj-い}
\vocabitem{面白い}{おもしろい}{interesante / divertido}{adj-い}
\vocabitem{難しい}{むずかしい}{difícil}{adj-い}
\vocabitem{易しい}{やさしい}{fácil}{adj-い}
\vocabitem{いい}{—}{bueno}{adj-い}
\bottomrule
\end{tabularx}
\end{tcolorbox}

\subsection*{🗣️ Frases Clave}

\frase{\ruby{京都}{きょうと}はとても\ruby{静}{しず}かです。}{Kioto es muy tranquilo.}
\frase{\ruby{日本語}{にほんご}の\ruby{勉強}{べんきょう}は\ruby{面白}{おもしろ}いです。}{El estudio de japonés es interesante.}
\frase{この\ruby{町}{まち}はあまり\ruby{賑}{にぎ}やかじゃありません。}{Este pueblo no es muy animado.}
\frase{パソコンはとても\ruby{便利}{べんり}ですね。}{Las computadoras son muy convenientes, ¿verdad?}

\subsection*{⚙️ Gramática}

\patron{Adjetivos-い vs Adjetivos-な}
  {Adj-い: termina en い (ej: 忙しい) \quad | \quad Adj-な: no termina en い o es una excepción (ej: きれい)}
  {\ej{この\ruby{本}{ほん}は\ruby{難}{むずか}しいです。}{Este libro es difícil.}
  \ej{この\ruby{町}{まち}は\ruby{賑}{にぎ}やかです。}{Este pueblo es animado.}}

\patron{Negación de adjetivos}
  {Adj-い: quitar い → くないです \quad | \quad Adj-な: adj + じゃありません}
  {\ej{\ruby{今日}{きょう}は\ruby{忙}{いそが}しくないです。}{Hoy no estoy ocupado.}
  \ej{\ruby{明日}{あした}は\ruby{暇}{ひま}じゃありません。}{Mañana no estoy libre.}}

\begin{tcolorbox}[ejerciciobox, title={📝 Ejercicios Rápidos}]
\begin{enumerate}
  \item Niega el adjetivo 面白い.
  \item Niega el adjetivo 便利.
  \item Traduce: \ruby{静}{しず}かな\ruby{公園}{こうえん}で\ruby{本}{ほん}を\ruby{読}{よ}みます。
\end{enumerate}
\vspace{0.4em}
\textbf{Respuestas:}
1) \ruby{面白}{おもしろ}くないです\quad
2) \ruby{便利}{べんり}じゃありません\quad
3) Leo un libro en un parque tranquilo.
\end{tcolorbox}

\begin{tcolorbox}[kanjibox, title={🀄 Kanjis de la Unidad}]
\begin{tabularx}{\linewidth}{cllXX}
\toprule
\textbf{Kanji} & \textbf{On} & \textbf{Kun} & \textbf{Significado} & \textbf{Vocab} \\
\midrule
\kanjirow{新}{しん}{あたら}{nuevo}{\ruby{新しい}{あたらしい} / \ruby{新聞}{しんぶん}}
\kanjirow{古}{こ}{ふる}{viejo}{\ruby{古い}{ふるい} / \ruby{中古}{ちゅうこ}}
\kanjirow{難}{なん}{むずか}{difícil}{\ruby{難しい}{むずかしい} / \ruby{困難}{こんなん}}
\kanjirow{便}{べん・びん}{たよ}{conveniente / noticia}{\ruby{便利}{べんり} / \ruby{便}{たよ}り}
\bottomrule
\end{tabularx}
\end{tcolorbox}
